\documentclass[11pt]{article}
\usepackage[a4paper,margin=28mm]{geometry}
\usepackage{amsmath,amssymb,amsthm,mathtools,hyperref}
\newcommand{\ddc}{dd^c}
\newcommand{\Var}{\operatorname{Var}}
\newtheorem{theorem}{Theorem}
\newtheorem{lemma}[theorem]{Lemma}
\theoremstyle{remark}\newtheorem{remark}[theorem]{Remark}
\title{Flat strips, half-planes and planes in nonclosed background cones}
\author{}
\date{}
\begin{document}
\maketitle
\begin{abstract}
We prove a complete metric-completion theorem for fixed three-dimensional
warped cones on a compact K\"ahler manifold times a curve. Their full
fixed-density Gauduchon volume is an explicit reciprocal integral, attained
by smooth balanced metrics. The affine Hessian metric on the unit-volume
surface is flat. Under two-sided finite-order tubular hypotheses at both
extreme level sets, the completion is a closed strip, a closed half-plane,
or the whole Euclidean plane. Every finite boundary point has a smooth
semipositive nef nonbig representative of volume one. A projective
threefold example gives a closed strip with explicit defining integral.
A flat complex torus gives an explicit complete Euclidean plane.
This is a specified-cone theorem, not an arbitrary-cone classification.
\end{abstract}

Use $\ddc=\frac{i}{2\pi}\partial\bar\partial$ and ordinary wedge powers.
All quasi-psh functions have their genuine canonical representatives.
For a smooth positive probability $\mu$ on an $n$-fold define
\[
 D_\mu(\gamma)=\int_X(\mu/\gamma^n)^{1/n}\gamma^n,
 \quad\lambda_\mu(\theta)=\inf_{\ddc\gamma^{n-1}=0}
 \frac{\int_X\theta\wedge\gamma^{n-1}}{D_\mu(\gamma)},
 \quad V_\mu=\lambda_\mu^n,
\]
where every smooth positive Gauduchon metric participates. On a smooth
strict positivity cone the metric is the affine Hessian
$g=-D^2\log\lambda_\mu$. Distances below are intrinsic to the specified
cone and its unit-volume hypersurface. At a nef boundary use decreasing
positive Hermitian regularization, equivalently the nonnegative full infimum.

\section{The entire warped cone}
Let $(B,\kappa)$ be a connected compact K\"ahler $m$-fold, $m\geq2$, with
$\int_B\kappa^m=1$. Let $C$ be a connected compact complex curve with
a smooth positive form $w$ of integral one. Fix a smooth strictly positive
probability volume $\sigma$ on $B$ and a smooth nonconstant function $a$
with $\min a=0$, $\max a=M>0$. Put $X=B\times C$, $n=m+1$,
$\mu=\sigma\wedge w$, and suppress pullbacks. In the three-dimensional
space spanned by $\kappa,aw,w$, write
\[
 \theta_{b,c,e}=b\kappa+Lw,\qquad L=ca+e.
\]
Its full cone of smooth strict potential representatives is
\begin{equation}\label{eq:cone}
 \mathcal C=\{b>0,\ e>0,\ e+cM>0\}.
\end{equation}
Indeed a smooth positive representative tested against the closed
semipositive form $\kappa^{m-1}\wedge w$ has integral $b>0$.
Its restriction to each compact curve fiber has positive area $L(y)$:
the integral of the restricted smooth $\ddc$ potential is zero. These
conditions are precisely \eqref{eq:cone}, and conversely $\theta$ itself
is positive there. The same tests applied to $\theta=\ddc u$ force
$b=0$ and $ca+e=0$ everywhere, hence $c=e=0$. Thus this space has no
nonzero smooth $\ddc$ direction. Only smooth potentials were restricted.
The backgrounds with $c\ne0$ are genuinely nonclosed.

\section{Exact volume and actual minimizing metrics}
Write $H(c,e)=\int_B\sigma/L$. Smooth Monge--Amp\`ere existence in the
closed positive base class \cite[Corollary 1]{TW} gives
\[
 A=b\kappa+\ddc\phi>0,\qquad A^m=\frac{b^m}{H}\frac{\sigma}{L}.
\]
The target has total mass $b^m$, so closed-class integration fixes the
source scalar constant. The source uses $i\partial\bar\partial$;
our potential is $2\pi$ times its potential. No endpoint equation is used.
Set
\[
 \Omega=A+Lw,\qquad\gamma=L^{-1/m}\Omega.
\]
Then
\begin{equation}\label{eq:cofactor}
 \Omega^n=\frac{nb^m}{H}\mu,
 \qquad\gamma^m=L^{-1}A^m+mA^{m-1}\wedge w.
\end{equation}
The first summand is a base top-degree form and the second is closed.
Consequently $\gamma$ is an actual smooth positive balanced test.
For every smooth positive Gauduchon metric $\eta$, smooth Stokes and
matrix AM--GM give
\[
 \int_X\theta\wedge\eta^m=\int_X\Omega\wedge\eta^m
 \geq(nb^m/H)^{1/n}D_\mu(\eta).
\]
For the named $\gamma$, the numerator is $nb^m$ and the denominator is
$(nb^m/H)^{m/n}H$. Hence equality is attained and
\begin{equation}\label{eq:volume}
 V_\mu(\theta)=\frac{nb^m}{H(c,e)}.
\end{equation}
This is the full infimum, with no mass law for nonclosed backgrounds.

The probability used below comes directly from this minimizing metric:
\begin{equation}\label{eq:trace}
 \frac{Lw\wedge\gamma^m}{b^m}=\frac{\sigma\wedge w}{HL}.
\end{equation}
Its base marginal is $\nu=\sigma/(HL)$. Thus the reciprocal-coefficient
moments controlling distance are normalized vertical-trace observables
of actual minimizing Gauduchon metrics. No pointwise endpoint limit of
these metrics is needed.

\section{The affine Hessian and flat coordinates}
Set
\[
 v=c/e\in(-1/M,\infty),\quad Q(v)=\int_B\frac{\sigma}{1+va},
 \quad d\nu_v=\frac{\sigma}{Q(v)(1+va)},\quad
 \ell_v=\frac{a}{1+va}.
\]
Equation~\eqref{eq:volume} gives
$g=\frac{m}{n}(db/b)^2+\frac1nD^2\log H$ in affine coordinates
$(b,c,e)$. If $F=\delta L/L$, then
\[
 D^2\log H(\delta,\delta)=2\mathbb E_\nu F^2-(\mathbb E_\nu F)^2.
\]
This is positive unless $\delta L=0$ everywhere. Nonconstancy of $a$
then forces $\delta c=\delta e=0$, proving positivity of the ambient metric.

On $\mathcal S=\{V_\mu=1\}$ put $U=\log b$. It is globally parameterized by
\[
 e=\frac{Q(v)}{n e^{mU}},\qquad c=ve.
\]
To pull back the affine Hessian correctly, let $t=de/e$, $z=Q'/Q$.
Along an affine variation $v''=-2t v'$, and $\log H=-\log e+\log Q(v)$
gives
\[
 D^2\log H=t^2-2zt\,dv+(Q''/Q-z^2)dv^2.
\]
On $\mathcal S$, $t=z\,dv-m\,dU$, so
\begin{equation}\label{eq:flat}
 g_{\mathcal S}=m\,dU^2+\rho(v)dv^2,
 \quad\rho(v)=\frac{Q''/Q-2(Q'/Q)^2}{n}
 =\frac2n\Var_{\nu_v}(\ell_v)>0.
\end{equation}
This is an affine-Hessian pullback, not the Hessian of the constant
restricted function. With $Y=\sqrt m\,U$ and
$R(v)=\int_{v_*}^v\sqrt{\rho(t)}\,dt$, it is a global isometry
onto $\mathbb R\times I$ for an open interval $I$.

\section{Both endpoint estimates}
Assume each of $\{a=0\}$ and $\{a=M\}$ is a finite disjoint union of
smooth compact real submanifolds of positive codimension. Near a component
$Z_i$ of the first set assume $a\asymp\operatorname{dist}(Z_i)^{k_i}$;
near a component of the second assume
$M-a\asymp\operatorname{dist}(Z_i)^{k_i}$, with $k_i>0$ and positive
two-sided constants. Off disjoint small tubes the corresponding function
is bounded below. Define separately
\[
 r_0=\min_{\{a=0\}}\frac{\operatorname{codim}_{\mathbb R}Z_i}{k_i},
 \qquad r_M=\min_{\{a=M\}}\frac{\operatorname{codim}_{\mathbb R}Z_i}{k_i}.
\]
These assumptions exclude intersecting strata and infinite-order zeros.

For $f=a$ or $M-a$ write
\[
 H_f=\int_B\frac{\sigma}{f+\epsilon},\quad
 J_f=\int_B\frac{\sigma}{(f+\epsilon)^2},\quad
 K_f=\int_B\frac{\sigma}{(f+\epsilon)^3},\quad
 d\nu_{f,\epsilon}=\frac{\sigma}{H_f(f+\epsilon)}.
\]
Both ends of \eqref{eq:flat} have the same invariant expression
\begin{equation}\label{eq:endmetric}
 \rho(v)dv^2=G_f(\epsilon)d\epsilon^2,
 \qquad G_f=\frac2n\left[\frac{K_f}{H_f}
             -\left(\frac{J_f}{H_f}\right)^2\right]
 =\frac2n\Var_{\nu_{f,\epsilon}}\frac1{f+\epsilon}.
\end{equation}
At $v\to\infty$ use $\epsilon=1/v$, $f=a$:
$\ell_v=\epsilon-\epsilon^2/(a+\epsilon)$ and
$dv=-\epsilon^{-2}d\epsilon$. At $v\to-1/M$ use
$\epsilon=-1/v-M$, $f=M-a$:
$\ell_v=(M+\epsilon)^2/(f+\epsilon)-(M+\epsilon)$ and
$dv=(M+\epsilon)^{-2}d\epsilon$. These give \eqref{eq:endmetric}
exactly. Its probabilities agree with the actual metric trace in
\eqref{eq:trace}, since positive scaling of $L$ leaves that marginal fixed.

Let $r$ be the corresponding minimum codimension/order ratio. Tubular
integration with smooth positive $\sigma$, splitting the normal radius
at $\epsilon^{1/k_i}$, gives
\[
 H_f\asymp\begin{cases}\epsilon^{r-1}&r<1,\\
 \log(1/\epsilon)&r=1,\\1&r>1,\end{cases}
 \qquad K_f\asymp\begin{cases}\epsilon^{r-3}&r<3,\\
 \log(1/\epsilon)&r=3,\\1&r>3.\end{cases}
\]
For $r>1$, $H_f$ tends to its finite positive endpoint integral. The
analogous threshold for $J_f$ is two. A finite sum of tube contributions
is governed by the smallest ratio.

For $r<1$ one has $G_f\asymp\epsilon^{-2}$. For the lower bound choose
a component realizing $r$ and two normal annuli of radii fixed multiples
of $\epsilon^{1/k_i}$. Choose those multiples far enough apart that the
two-sided bounds separate their $f/\epsilon$ ranges. Each annulus has
$\nu_{f,\epsilon}$ probability bounded below, since its weighted mass
and $H_f$ both have order $\epsilon^{r-1}$. The reciprocal values differ
by at least a fixed multiple of $\epsilon^{-1}$. Apply
$\Var X=\frac12\iint(X(x)-X(y))^2d\nu(x)d\nu(y)$.
The upper bound is $X\leq\epsilon^{-1}$.

For $r=1$ the near annulus has probability at least
$c/\log(1/\epsilon)$. In the same tube, normal radii between
$\epsilon^{1/(2k_i)}$ and a fixed small radius have probability bounded
below independently of $\epsilon$, by the logarithmic integral.
Reciprocal values there are at most $C\epsilon^{-1/2}$, whereas on the
near annulus they are at least $c\epsilon^{-1}$. The same variance
identity and $G_f\leq2K_f/(nH_f)$ yield
\[
 G_f\asymp\frac1{\epsilon^2\log(1/\epsilon)}.
\]
Both cases have infinite length. If $r>1$, the bound
$G_f\leq C K_f$ makes $\sqrt{G_f}$ integrable: the worst power is
$\epsilon^{(r-3)/2}$ for $1<r<3$, with exponent greater than $-1$;
at three it is $\sqrt{\log(1/\epsilon)}$, and above three it is bounded.
Thus an end is finite exactly when $r>1$. These quantitative estimates
are on the minimizing metrics' trace moments, not a formal distance test.

\section{Every Cauchy boundary point}
\begin{theorem}\label{th:completion}
Under the stated hypotheses the metric completion of $\mathcal S$ is
a closed Euclidean strip if $r_0,r_M>1$, a closed Euclidean half-plane
if exactly one exceeds one, and the whole Euclidean plane otherwise.
Every finite end is parameterized by $b>0$ and has the smooth
semipositive nef nonbig representatives of volume one displayed below.
There are no other Cauchy boundary points in this fixed cone.
\end{theorem}
\begin{proof}
The isometry onto $\mathbb R\times I$ identifies its intrinsic distance
with the Euclidean norm: the interval product is convex and its straight
segments are admissible. Its completion is its Euclidean closure. A
Cauchy sequence has finite $Y$, hence a finite positive limit of $b$;
its $R$ coordinate can reach only an interior point or a finite end.
The endpoint estimates give the stated alternatives.

If $r_0>1$ put $H_a=\int_B\sigma/a$. Since $vQ(v)\to H_a$, the
coefficient limits at the right end are exactly
\[
 \theta_+(b)=b\kappa+\frac{H_a}{nb^m}a w.
\]
If $r_M>1$ put $H_-=\int_B\sigma/(1-a/M)$. At the left end they are
\[
 \theta_-(b)=b\kappa+\frac{H_-}{nb^m}(1-a/M)w.
\]
These forms are smooth semipositive, nef by positive Hermitian
perturbation, and nonclosed. Increasing the vertical coefficient by
$\epsilon$ and using \eqref{eq:volume} gives a limit of one. For every
fixed admissible metric the quotient converges to its endpoint quotient.
Monotonicity gives the opposite inequality for the infima, so the full
endpoint infimum is that limit. The same fixed-test argument proves
agreement with the original Hermitian regularization. No singular
endpoint equation or undefined top product is required.

To exclude bigness for every genuine quasi-psh $u$, suppose
$\theta_\pm+\ddc u\geq\delta(\kappa+w)$. Test against the pullback of
any smooth nonnegative base top-degree form $\eta$. It is $\ddc$-closed
by dimension, so distributional Stokes gives
$\int_B(L_\pm-\delta)\eta\geq0$. The coefficient $L_\pm$ vanishes
somewhere; a nonzero $\eta$ supported where $L_\pm<\delta$ contradicts
this inequality. No singular potential or current is restricted to a fiber.
\end{proof}

\section{A projective closed strip}
Take $B=\mathbb P^2$, $\sigma=\kappa^2$ for normalized Fubini--Study,
$C=\mathbb P^1$, and
\[
 a([Z])=\frac{|Z_1|^2/2+|Z_2|^2}{|Z_0|^2+|Z_1|^2+|Z_2|^2}.
\]
Here $M=1$. Both extreme sets are single points, at $[1:0:0]$ and
$[0:0:1]$, with positive quadratic behavior in their affine charts.
Both real codimensions are four and orders two, so $r_0=r_M=2$.

For completeness the law of $a$ under $\kappa^2$ has density $4s$ on
$[0,1/2]$ and $4(1-s)$ on $[1/2,1]$. In the $Z_0$ chart, polar
integration gives $2\,ds_1ds_2/(1+s_1+s_2)^3$ for $s_j=|z_j|^2$.
The transformation $u_j=s_j/(1+s_1+s_2)$ yields uniform density two
on $u_1,u_2\geq0$, $u_1+u_2\leq1$. The sublevel triangle for
$a=u_1/2+u_2\leq s\leq1/2$ has area $s^2$, giving density $4s$.
Exchanging $u_0,u_2$ gives the other half. Consequently
\[
 H_a=H_-=4\log2,\qquad
 Q(v)=\frac4{v^2}\big[(v+1)\log(1+v)-(v+2)\log(1+v/2)\big]
\]
for $v\ne0$, $v\in(-1,\infty)$, with smooth removable value $Q(0)=1$.
The formula follows by integrating the two linear densities; the integral
definition verifies the extension at zero. The cone is exactly
$b>0,e>0,e+c>0$, and its two finite faces are
\[
 b\kappa+\frac{4\log2}{3b^2}a w,
 \qquad b\kappa+\frac{4\log2}{3b^2}(1-a)w.
\]
They are nef nonbig and have full volume one. Its exact strip width is
the positive finite integral
\[
 W=\int_{-1}^{\infty}
 \sqrt{\frac{Q''/Q-2(Q'/Q)^2}{3}}\,dv.
\]
The coordinates $Y=\sqrt2\log b$ and $R(v)$ give the ordinary Euclidean
distance, also on the completion.

\section{An explicit complete plane on a torus}
Let $E=\mathbb C/(\mathbb Z+i\mathbb Z)$, $B=E^m$, $m\geq2$,
and write $z_j=u_j+i v_j$ on its cover. Put
\[
 \eta_j=\tfrac i2 dz_j\wedge d\bar z_j=du_j\wedge dv_j,
 \quad s=(m!)^{-1/m},\quad\kappa=s\sum_j\eta_j,
 \quad\sigma=\kappa^m=\prod_j\eta_j.
\]
These descend and give a normalized flat K\"ahler form and probability.
Take $a=\sin^2(\pi u_1)$, a globally smooth function with $M=1$.
The minimum and maximum sets are the real hypersurfaces $u_1=0$ and
$u_1=1/2$ modulo one. Both have quadratic normal behavior, so both
codimension/order ratios are $1/2$. For $v=c/e>-1$,
\[
 Q(v)=\int_0^1\frac{du}{1+v\sin^2(\pi u)}=\frac1{\sqrt{1+v}}.
\]
Splitting the period at $1/2$ and substituting $t=\tan(\pi u)$ gives
$\frac2\pi\int_0^\infty[1+(1+v)t^2]^{-1}dt$, proving the formula.

This example also admits an elementary explicit base construction. Set
$F_v(u)=[Q(v)(1+v\sin^2(\pi u))]^{-1}$, whose period integral is one.
For each $b>0$, twice integrating the zero-mean smooth periodic right side
gives a smooth periodic $\phi_v$ (with its $b$ dependence suppressed)
with $\phi_v''=4\pi bs(F_v-1)$: choose the
first integration constant so that its derivative also has mean zero.
Our convention gives $\ddc\phi_v=\phi_v''\eta_1/(4\pi)$, hence
\[
 A=b\kappa+\ddc\phi_v
 =bs\left(F_v\eta_1+\sum_{j=2}^m\eta_j\right)>0,
 \qquad A^m=b^mF_v\sigma=\frac{b^m}{H}\frac{\sigma}{L}.
\]
Thus the actual minimizing metric and full-test equality follow directly
from \eqref{eq:cofactor}, without an existence theorem for this example.
The entire cone is $b>0,e>0,e+c>0$. Differentiating its explicit $Q$ yields
\[
 \rho(v)=\frac1{4n(1+v)^2},\qquad
 g_{\mathcal S}=m\,dU^2+\frac1{4n}\left(\frac{dv}{1+v}\right)^2.
\]
The exact coordinates $Y=\sqrt m\log b$ and
$R=\log(1+v)/(2\sqrt n)$ range over the entire Euclidean plane.
Consequently this unit-volume surface is already complete and has no
finite Cauchy boundary, although its coefficient cone has two vertical faces.
Its backgrounds with $c\ne0$ are genuinely nonclosed.

\begin{remark}
The fixed cone and product probability are part of the theorem. We do
not classify arbitrary nonclosed cones, intersecting extreme strata or
infinite-order zeros. No uniqueness or pointwise convergence of minimizing
metrics, endpoint Monge--Amp\`ere solvability, or literature priority is asserted.
\end{remark}
\begin{thebibliography}{9}
\bibitem{TW} V. Tosatti and B. Weinkove, \emph{The complex Monge--Amp\`ere
equation on compact Hermitian manifolds}, arXiv:0910.1390v2 (11 November 2009),
Corollary 1 and closed-background normalization, pp.~1--2;
\url{https://arxiv.org/abs/0910.1390}.
\end{thebibliography}
\end{document}
